Unfolding \(\mathsf W_j\) in \(\text{def:terminal-qe-test}\) gives
\[
\mathsf C_{\mathcal G}(\theta^i;P_{\mathrm{obs}})
\quad(i=0,1),
\qquad
q_{a,y}(\theta^0)\neq q_{a,y}(\theta^1).
\tag{1}
\]
Apply the converse direction of \(\text{thm:response-fiber-equivalence}\) separately to the two compatibility conjuncts.  It decodes \(\theta^i\) to a model \(\mathcal M_i\) satisfying
\[
\mathcal M_i\in\mathfrak F(\mathcal G,P_{\mathrm{obs}}),
\qquad
Q_a^{\mathcal M_i}(y)=q_{a,y}(\theta^i),
\quad i=0,1.
\tag{2}
\]
The decoding uses the unique one-hot response-table selector and the product measure \(\prod_h\rho_h^i\), so it retains one mutually independent factor for each individual bidirected edge and each private vertex factor.  The equalities \(p_o(\theta^i)=P_{\mathrm{obs}}(o)\) in \(\mathsf C_{\mathcal G}\), together with \(\text{def:observational-fiber}\), give
\[
P_{\mathcal M_0}(\mathcal O)=P_{\mathcal M_1}(\mathcal O)=P_{\mathrm{obs}}.
\tag{3}
\]
Likewise, the strict inequalities \(p_{v,r}(\theta^i)>0\) in \(\mathsf C_{\mathcal G}\), and the full-cell equality supplied by \(\text{thm:response-fiber-equivalence}\), prove strict positivity of every full-data cell of both decoded models.  Finally, substituting (2) into the last inequality in (1) yields
\[
Q_a^{\mathcal M_0}(y)\neq Q_a^{\mathcal M_1}(y).
\tag{4}
\]
Equations (2)--(4) prove all conclusions.  The exact-specialization premise is used to supply the represented points \(\theta^0,\theta^1\), not to establish the decoding identities.